%! Solving Quadratics by Factoring — Unit 3, Lesson 3.2 — Lesson Plan
\documentclass[10pt]{article}
\usepackage{algebra2-article}
\usepackage{algebra2-boxes}
\usepackage{graphicx}   % -article does NOT load graphicx; needed for thumbnails/screenshots
\graphicspath{{images/}}
\newcommand{\TallMath}[1]{$\displaystyle #1\rule[-1.4em]{0pt}{3.2em}$}
\newcommand{\CourseName}{Algebra 2: Shepherd}
\newcommand{\SchoolYear}{2026--2027}
\newcommand{\MeetingLength}{55 minutes}
\newcommand{\UnitNumberName}{Unit 3: Quadratic Functions \quad}
\newcommand{\LessonNumberName}{Lesson 3.2: Solving Quadratics by Factoring}

\begin{document}
\begin{center}
    {\Large\bfseries \CourseName: \SchoolYear \\
    {\small\bfseries \UnitNumberName \LessonNumberName}}
\end{center}

\begin{tcolorbox}[colback=forestbg, colframe=forest, boxrule=0.9pt, arc=2mm,
    left=2mm, right=2mm, top=1.5mm, bottom=1.5mm]
    \textbf{Primary Objective:} Students can \emph{produce} the factored form of a quadratic and use it to
    solve. They factor a quadratic completely --- pulling out a \emph{GCF}, factoring $x^2+bx+c$ by finding
    two integers with the right product and sum, factoring $ax^2+bx+c$ by grouping, and recognizing the
    \emph{difference of squares} and \emph{perfect-square trinomial} patterns --- then apply the
    \emph{Zero Product Property} to solve $ax^2+bx+c=0$. They connect the solutions (the \emph{roots}) to
    the parabola's \emph{$x$-intercepts / zeros}: where Lesson~3.1 \emph{read} the zeros off a given
    factored form, Lesson~3.2 lets them \emph{find} that form by hand and read the zeros from any
    factorable standard-form quadratic.
    \\[2pt]
    \textbf{Standards (2023 VA SOL):} \textbf{A2.EO.3b} (factor polynomials completely in one and two
    variables over the integers) and \textbf{A2.EO.3d} (recognize and use the difference-of-squares and
    perfect-square-trinomial identities); \textbf{A2.EI.2b} (solve a quadratic equation in one variable
    algebraically --- here, the factoring method, restricted to real, factor-over-integers cases;
    square-root/completing-the-square follow in 3.3 and complex roots in 3.4--3.5), \textbf{A2.EI.2a}
    (create a quadratic equation to model a context), and \textbf{A2.EI.2d} (verify solutions graphically
    and interpret them --- the roots are the parabola's $x$-intercepts).
\end{tcolorbox}

\renewcommand{\arraystretch}{1.4}
\begin{skillbox}[Priority Ideas \& Skills]{goldbox}
\begin{tabularx}{\linewidth}{@{}X|X@{}}
\textbf{Priority Ideas \& Skills} & \textbf{Key Understandings (the ``why'')} \\[2pt]
\hline
\vspace{-6pt}
\begin{itemize}[leftmargin=1.1em, nosep, after=\vspace{2pt}]
  \item Always factor out the \textbf{GCF} first, then factor what remains completely.
  \item Factor $x^2+bx+c$: find two integers whose \textbf{product is $c$} and \textbf{sum is $b$}.
  \item Factor $ax^2+bx+c$ ($a\ne1$) by \textbf{grouping} (the $a\!\cdot\!c$ method).
  \item Recognize \textbf{difference of squares} $x^2-k^2$ and \textbf{perfect-square trinomials}.
  \item Solve by the \textbf{Zero Product Property}: set $=0$, factor, set each factor to $0$.
  \item Connect the \textbf{roots} to the parabola's \textbf{$x$-intercepts / zeros}.
\end{itemize}
&
\vspace{-6pt}
\begin{itemize}[leftmargin=1.1em, nosep, after=\vspace{2pt}]
  \item Factoring is \emph{un-multiplying}: it reverses the FOIL that built the trinomial.
  \item The product/sum search works because $(x+p)(x+q)=x^2+(p+q)x+pq$.
  \item A product equals zero \emph{only} when a factor is zero --- that is what turns a factored equation into simple linear ones.
  \item You must move everything to one side (\,$=0$) \emph{before} the Zero Product Property applies --- it is a fact about zero, not any number.
  \item The roots are exactly where the graph crosses the $x$-axis, tying algebra back to the picture from 3.0--3.1.
\end{itemize}
\end{tabularx}
\end{skillbox}

\begin{skillbox}[{Vocabulary, Concepts, \& Theorems}]{sky}
\renewcommand{\arraystretch}{1.3}
\begin{tabularx}{\linewidth}{@{}l X@{}}
\textbf{Factoring} & Writing a polynomial as a product of factors --- the reverse of multiplying/expanding. \\
\textbf{GCF} & Greatest common factor; the largest monomial dividing every term, factored out first. \\
\textbf{Zero Product Property} & If $AB=0$, then $A=0$ or $B=0$ (or both). The engine of solving by factoring. \\
\textbf{Roots / zeros / solutions} & Three names for the $x$-values that make $f(x)=0$ --- the parabola's $x$-intercepts. \\
\textbf{Difference of squares} & $x^2-k^2=(x-k)(x+k)$. \\
\textbf{Perfect-square trinomial} & $x^2\pm2kx+k^2=(x\pm k)^2$; factoring it gives one repeated (double) root. \\
\textbf{Double root} & A root that occurs twice; the parabola \emph{touches} the $x$-axis there (vertex on the axis). \\
\end{tabularx}
\end{skillbox}

\begin{fixedskillbox}[Activate Prior Knowledge \& Spiral Review]{forestbg}
    The Warm-Up rehearses exactly what factoring needs: \textbf{(1)} \emph{multiplying} binomials with FOIL
    (the operation factoring will reverse); \textbf{(2)} the number game at the heart of it --- finding two
    integers with a given \emph{product} and \emph{sum}, including the sign cases; and \textbf{(3)} reading
    a parabola's \emph{$x$-intercepts} off a pre-drawn graph and solving the trivial linear equations
    $x+p=0$, so the Zero Product Property lands as ``two easy equations.'' Debrief item~1 next to item~2:
    factoring is item~1 run backward.
\end{fixedskillbox}

\begin{skillbox}[Hook]{forestbg}
    Put the unit's anchor parabola back on the board:
    \[ f(x)=x^2-2x-3. \]
    In Lesson~3.1 you were \emph{handed} its factored form $f(x)=(x+1)(x-3)$ and read the zeros $-1$ and
    $3$ for free. Ask: \textbf{where did that factored form come from?} Today we manufacture it. Once we
    have $(x+1)(x-3)$, finding where the graph crosses the $x$-axis is a two-line problem: a product is
    zero only when one of its pieces is zero, so $x+1=0$ or $x-3=0$. Land the theme: \emph{factoring turns
    one hard quadratic equation into two easy linear ones --- and hands us the $x$-intercepts by hand.}
\end{skillbox}

\begin{skillbox}[Lesson]{forestbg}
\begin{multicols}{2}
\textbf{1. The Zero Product Property.} Start with the idea, not the algebra: if two numbers multiply to
$0$, at least one \emph{is} $0$ --- nothing else works. So $(x+1)(x-3)=0$ splits into $x+1=0$ or
$x-3=0$, giving $x=-1$ or $x=3$. Show those are the $x$-intercepts of the anchor parabola. Stress the
precondition: the product must equal \emph{zero}, so everything moves to one side first.

\textbf{2. Factor $x^2+bx+c$.} Un-FOIL: since $(x+p)(x+q)=x^2+(p+q)x+pq$, we need two integers whose
\emph{product} is $c$ and \emph{sum} is $b$. Work $x^2-2x-3$: product $-3$, sum $-2$ $\Rightarrow$ $+1$
and $-3$, so $(x+1)(x-3)$ --- the anchor, produced by hand. Drill the sign logic: product negative
$\Rightarrow$ opposite signs; product positive $\Rightarrow$ same sign as $b$.

\textbf{3. GCF first.} Before anything, pull out the greatest common factor: $2x^2-8x=2x(x-4)$, and
solving $2x^2-8x=0$ gives $x=0$ or $x=4$. Use this to flag the \#1 error --- \emph{never divide both
sides by $x$}; that silently deletes the root $x=0$.

\columnbreak

\textbf{4. Factor $ax^2+bx+c$ ($a\ne1$) by grouping.} The $a\!\cdot\!c$ method: for $2x^2+7x+3$,
$a\!\cdot\!c=6$; find two numbers with product $6$, sum $7$ ($6$ and $1$); split the middle term
$2x^2+6x+x+3$; group $2x(x+3)+1(x+3)=(2x+1)(x+3)$. Roots $x=-\tfrac12,\,-3$.

\textbf{5. Special patterns.} \emph{Difference of squares} $x^2-16=(x-4)(x+4)$ (roots $\pm4$);
\emph{perfect-square trinomial} $x^2-6x+9=(x-3)^2$ (one \textbf{double root} $x=3$). Tie the double root
to the graph: the parabola's vertex sits \emph{on} the $x$-axis, so it touches rather than crosses.

\textbf{6. Solve and verify.} The full routine: (i) set $=0$, (ii) factor completely, (iii) apply the
Zero Product Property, (iv) solve each linear piece, (v) check the roots are the graph's $x$-intercepts.
Preview: Lesson~3.3 handles the quadratics that \emph{won't} factor over the integers --- square roots
and completing the square.
\end{multicols}
\end{skillbox}

\begin{skillbox}[Explicit Instruction: The Factor-and-Solve Routine]{forestbg}
\begin{tabularx}{\linewidth}{@{}p{0.46\linewidth} X@{}}
\textbf{Steps (post and refer back):}
\begin{enumerate}[leftmargin=1.4em, nosep, itemsep=2pt]
  \item Write the equation with $0$ on one side.
  \item Factor out the GCF.
  \item Factor what remains completely (product/sum, grouping, or a special pattern).
  \item Set \emph{each} factor equal to $0$.
  \item Solve the linear equations; list all roots.
  \item Verify: substitute back, or check the roots are the $x$-intercepts.
\end{enumerate}
&
\textbf{Worked example.} Solve $x^2-x-6=0$.
\begin{itemize}[leftmargin=1.2em, nosep, itemsep=1pt]
  \item Already $=0$; no GCF.
  \item Product $-6$, sum $-1$ $\Rightarrow$ $-3,\,+2$: $(x-3)(x+2)=0$.
  \item $x-3=0$ or $x+2=0$.
  \item $x=3$ or $x=-2$.
  \item Check: $3^2-3-6=0$ \checkmark; $(-2)^2-(-2)-6=0$ \checkmark. These are the $x$-intercepts.
\end{itemize}
\end{tabularx}
\end{skillbox}

\begin{skillbox}[Active Monitoring]{redbox}
Circulate during the notes practice and Tier work. Watch for and correct:
\begin{itemize}[leftmargin=1.2em, nosep]
  \item \textbf{dividing by $x$}: solving $3x^2=12x$ by cancelling an $x$ and losing the root $x=0$ --- always move to one side and factor;
  \item applying the Zero Product Property when the equation is \emph{not} set to zero (e.g.\ $(x-2)(x-5)=4$ does \emph{not} give $x-2=4$);
  \item \textbf{sign errors} in the product/sum search --- forgetting that a negative $c$ forces opposite signs;
  \item skipping the \textbf{GCF}, then being unable to factor what remains;
  \item forgetting the \textbf{double root} counts once as a value but touches the axis (perfect-square case).
\end{itemize}
Cold-call prompts: ``Why can we split $(x+1)(x-3)=0$ into two equations --- but not $(x+1)(x-3)=5$?''
``Two integers multiply to $-12$ and add to $1$; what are they?'' ``Where are the roots on the graph?''
\end{skillbox}

\begin{skillbox}[Group Work \& Differentiation]{redbox}
\textbf{Zeros Factory} activity (tiered; mirrors the notes).
\begin{multicols}{3}
\textbf{Tier R --- Remediate.} Solve \emph{already-factored} equations with the Zero Product Property,
factor simple $x^2+bx+c$, and match each to the pre-drawn parabola whose $x$-intercepts it names.

\columnbreak
\textbf{Tier A --- Approaching.} Factor and solve a mix ($x^2+bx+c$ and GCF cases), complete a
\emph{equation $\to$ factored $\to$ roots} table, and confirm the roots against a labeled graph.

\columnbreak
\textbf{Tier E --- Extension.} Factor $ax^2+bx+c$ by grouping, use the difference-of-squares and
perfect-square patterns (spotting a double root), and \emph{model}: write and solve a quadratic for a
projectile/area context.
\end{multicols}
\end{skillbox}

\begin{skillbox}[Individual Work \& Assessment]{redbox}
\textbf{Exit Ticket} (independent, no notes): solve $x^2+2x-8=0$ by factoring and state the two zeros;
solve $3x^2-12x=0$ (GCF --- watch the $x=0$ root) and say which points are the graph's $x$-intercepts;
and an \textbf{SOL-style multiple-choice} item that turns on the sign flip between roots and factors
(zeros $-3,4 \Rightarrow (x+3)(x-4)$). Collect and sort into ``got it / GCF slip (lost $x=0$) /
sign-flip slip / Zero-Product-without-zero slip'' to decide whether Lesson~3.3 opens with a re-teach.
\end{skillbox}

\begin{skillbox}[Reinforcement \& Extension]{goldbox}
\textbf{Homework:} multiply then factor (see the reverse), factor across all cases (GCF, $x^2+bx+c$,
$a\ne1$, difference of squares, perfect square), solve by factoring, and match roots to a pre-drawn
graph. \textbf{Extension:} a garden-border \emph{area} model --- write a quadratic, solve by factoring,
and reject the impossible (negative) root in context. \textbf{Preview:} Lesson~3.3, \emph{Solving by
Square Roots \& Completing the Square} --- what to do when a quadratic will not factor over the integers.
\end{skillbox}
\end{document}
